\documentclass[11pt]{article}
\usepackage[letterpaper,margin=1.02in,top=0.92in,bottom=0.92in]{geometry}
\usepackage{amsmath,amssymb,amsthm,booktabs,xcolor,fancyhdr}
\usepackage[colorlinks=true,allcolors=teal]{hyperref}
\definecolor{ink}{RGB}{25,63,80}
\newtheorem{theorem}{Theorem}
\newtheorem{lemma}[theorem]{Lemma}
\newtheorem{corollary}[theorem]{Corollary}
\newcommand{\F}{\mathbb F_q}
\newcommand{\Alt}{\operatorname{Alt}}
\newcommand{\Sp}{\operatorname{Sp}}
\newcommand{\TV}{d_{\rm TV}}
\newcommand{\rk}{\operatorname{rank}}
\newcommand{\nil}{\operatorname{null}}
\newcommand{\Prob}{\mathbb P}
\newcommand{\E}{\mathbb E}
\newcommand{\T}{\mathsf T}
\newcommand{\src}{https://github.com/VladimirReshetnikov/ProveIt/blob/35e2a5a634924e57ac6b77bd73a28570dd1a362b/Analysis/GeneralizedRiemannHypothesis/Article/sections/spn.tex}
\pagestyle{fancy}\fancyhf{}\lhead{\small\color{ink}Symplectic compressions}\rhead{\small\color{ink}Optimal tests and limiting distances}\cfoot{\thepage}
\setlength{\headheight}{14pt}\setlength{\parskip}{2.5pt}\raggedbottom
\hypersetup{pdftitle={Optimal tests for nondegenerate symplectic compressions},pdfauthor={Research note}}
\begin{document}
\begin{center}
{\large\color{ink}\sffamily FINITE GEOMETRY AND PROBABILITY}\par\vspace{8pt}
{\LARGE\bfseries Optimal tests for nondegenerate\\[3pt] symplectic compressions}\par\vspace{7pt}
{\small November 27, 2026}
\end{center}

\begin{abstract}
Given an observed rectangular block, how well can one distinguish a uniform symplectic transformation from a uniform matrix over the same finite field? For two nondegenerate frames of any dimensions, one threshold on the rank of a residual alternating form attains the exact total variation distance. Its meaning changes when the sum of the frame dimensions exceeds the ambient dimension. We compute the limiting distance for every integer offset from this boundary, with an error uniform in the offset and the aspect ratio. This includes complementary frames of fixed unequal dimension gap, an explicit open case in the ProveIt article. The proof uses a monotonicity property of alternating Gram counts and Fourier averaging over the orbit of a nonsingular form. All statements here are unconditional.
\end{abstract}

\section{The statistical experiment}
Fix a finite field $\F$, including characteristic two. An alternating form has $B(x,x)=0$ for every $x$. In a symplectic space $(V,J)$ of dimension $2n$ choose two matrices $Y,X_0$ with respectively $a=2k$ and $b=2h$ linearly independent columns spanning nondegenerate subspaces, where $k,h>0$. Write $H_1=Y^\T JY$, $H_0=X_0^\T JX_0$. Draw $g$ uniformly from $\Sp(V)$ and compare
\[
 M=Y^\T JgX_0\qquad\text{with}\qquad
 Z\text{ uniform on }\F^{a\times b}.
\]
The coordinate block of $gX_0$ in the target has the same statistical distance, since it is $H_1^{-1}M$. No assumption about the intersection of the original frames is needed. For a candidate matrix $z$ define
\[
 S(z)=H_0+z^\T H_1^{-1}z,\qquad 2c=2n-a,\qquad s=c-h=n-k-h .
\]
Thus $2s$ measures how much ambient dimension is left after counting both frames. The case $s=0$ includes every complementary size. Write $D=\TV(M,Z)=\sup_E|\Prob(M\in E)-\Prob(Z\in E)|$. With equal prior probabilities the optimal success probability is $(1+D)/2$.

We use the products $\eta_j=\prod_{i=1}^j(1-q^{-2i})$, $\theta_j=\prod_{i=1}^j(1-q^{1-2i})$, and $\gamma_j=\prod_{i=1}^j(1-q^{-i})$. Empty products are one. Infinite subscripts mean convergent positive infinite products.

\section{One optimal decision rule in each regime}

\begin{theorem}[Exact test]\label{th:finite}
If $s\ge0$, the observed matrix favors the symplectic model exactly when $S(z)$ is singular. In this case
\begin{equation}
 D=(1-\alpha_s)\Prob\bigl(S(Z)\text{ is nonsingular}\bigr),
 \qquad \alpha_s=\frac{\eta_{h+s}\eta_{k+s}}{\eta_s\eta_{h+k+s}}.
 \label{eq:positive}
\end{equation}
If $s=-r<0$, every matrix with $\nil S(z)\ge 2r$ favors the symplectic model, and every other matrix has zero symplectic probability. In particular,
\begin{equation}
 D=\Prob\bigl(\nil S(Z)<2r\bigr).
 \label{eq:negative}
\end{equation}
Here $\nil$ means the dimension of the kernel.
\end{theorem}

\begin{theorem}[Uniform limiting profile]\label{th:limit}
For fixed $q$, all admissible dimensions satisfy
\begin{equation}
 |D-\mathcal D_q(s)|\ \le\ C_q(q^{-a}+q^{-b}),\qquad
 \mathcal D_q(s)=
 \begin{cases}
 \theta_\infty(1-\eta_\infty/\eta_s),&s\ge0,\\[2pt]
 \displaystyle\sum_{j=0}^{-s-1}
 \frac{\theta_\infty q^{-j(2j-1)}}{\gamma_{2j}},&s<0 .
 \end{cases}
 \label{eq:profile}
\end{equation}
The constant depends only on $q$. In particular, as $a,b\to\infty$, the profile tends to zero for $s\to+\infty$ and to one for $s\to-\infty$. At $s=0$ the limit is $\theta_\infty(1-\eta_\infty)$ for every aspect ratio.
\end{theorem}

For instance, over $\mathbb F_2$, the limiting distances at $s=-1,0,1$ are about $0.41942$, $0.13063$, and $0.03437$, respectively. A change of two in the ambient dimension can thus survive the limit in which both observed dimensions grow without bound.

\paragraph{Relation to the source.}
The ProveIt article \cite{proveit} derives a residual likelihood, gives the sign test at $s=0$, and computes the limiting full-matrix distance for square and increasingly tall complementary blocks. Its question \texttt{spn:q:symplectic} singles out the fixed positive gap $a-b$ at $s=0$. Theorem~\ref{th:limit} covers that gap and allows the offset $s$ itself to vary. We include proofs of the likelihood and the needed counting identities. The two new ingredients are the off-boundary sign test and a Fourier estimate for the residual corank which is uniform in the dimension of the alternating forms. The geometry of this model is related to existing counts of pairs of nondegenerate subspaces \cite{anzahl}. In the complementary case their graph-subspace counts already encode the scalar nonsingularity probability in \eqref{eq:positive}.

\section{Gram counts and monotonicity}

Write $J_v$ for a fixed standard symplectic form on $\F^{2v}$. For $T$ uniform in $\F^{2v\times b}$ put $G_v=T^\T J_vT$. Let $m=\binom b2$, and define the density $F_v(B)=q^m\Prob(G_v=B)$ relative to the uniform law on $\Alt_b$. Congruence makes $F_v$ a function of the rank of $B$. That rank is even in every characteristic.

\begin{lemma}[Residual likelihood]\label{lem:like}
The likelihood of $z$ under the symplectic model, relative to uniform, is $F_c(S(z))/\Pi$, where $\Pi=\eta_n/\eta_{n-h}$. Its support is $\rk S(z)\le2c$.
\end{lemma}
\begin{proof}
A frame with Gram matrix $H_0$ extends to a symplectic basis after a fixed choice of basis on its domain. Thus $gX_0$ is uniform on these frames. A uniform $2n\times b$ matrix conditioned to have Gram matrix $H_0$ has the same law, since that condition implies injectivity. Decompose it into target coordinates and independent orthogonal-complement coordinates. Specifying the observed matrix $z$ fixes its target component $YH_1^{-1}z$, with Gram matrix $-z^\T H_1^{-1}z$. Its other component must have Gram matrix $S(z)$. The likelihood is consequently $F_c(S(z))/F_n(H_0)$.

In a symplectic space of dimension $2v$, the number of ordered pairs with pairing one is $(q^{2v}-1)q^{2v-1}$. Successively splitting off such pairs, a given nonsingular $b$-form has $\prod_{i=0}^{h-1}(q^{2(v-i)}-1)q^{2(v-i)-1}$ realizations for $v\ge h$. Multiplication by $q^{m-2vb}$ gives $F_v(H_0)=\eta_v/\eta_{v-h}$. For the support, a form of rank $2j\le2c$ can be realized by embedding its nondegenerate quotient in the complement and killing its radical. The opposite inequality precludes a realization.
\end{proof}

\begin{lemma}[Monotonicity of the point probabilities]\label{lem:mono}
For all $v\ge0$, $F_v(B)$ is nonincreasing with $\rk B$.
\end{lemma}
\begin{proof}
Regard an alternating matrix as a bivector. Passing from $G_v$ to $G_{v+1}$ adds $u\wedge w=uw^\T-wu^\T$ for independent uniform column vectors $u,w\in\F^b$. This increment is symmetric in law under negation. The transition operator on densities therefore agrees with its operator on functions. On half-ranks it is a nearest-neighbor chain, since one step changes rank by at most two.

At rank $2i$ write $L$ for the column space, and $e=b-2i$. An increase occurs exactly when $u,w$ project to independent vectors modulo $L$. Indeed in that case a complement containing $u,w$ exhibits the two extra ranks in a separate block. In the other case the new column space lies in a space of dimension at most $2i+1$. Thus the upward probability is
\[
 p_i^+=(1-q^{-e})(1-q^{1-e})\qquad(e\ge2).
\]
A decrease from the next rank $2i+2$ requires both vectors in its column space. For a nonzero decomposable increment this follows from equality in
$\rk B\le\rk(B+u\wedge w)+\rk(u\wedge w)$, which forces its column space into that of $B$. On the latter space the block of $B$ is invertible. Given a nonzero first vector, singularity after the update prescribes one nonzero value of the linear functional $w\mapsto w^\T B^{-1}u$ in these block coordinates. Indeed the determinant of the updated block divided by $\det B$ is $(1+w^\T B^{-1}u)^2$. This identity holds also in characteristic two. The downward probability is consequently exactly
\[
 p_{i+1}^-=q^{-2(e-2)-1}(1-q^{-2(i+1)})\le q^{1-e}\le 1-p_i^+.
\]
For two adjacent starting states the last inequality orders the probability of crossing their common boundary. At every other threshold the ordering is automatic. Thus the transition preserves nonincreasing functions of half-rank. The initial density, supported at rank zero, is nonincreasing. Iteration proves the claim.
\end{proof}

\paragraph{An exact calculation with $h+1$ states.}
The preceding proof gives an inexpensive evaluation of both finite distances. Start at half-rank $h$, the rank state of $H_0$. At state $i$ put $e=b-2i$ and take the transitions
\[
 p_i^+=(1-q^{-e})(1-q^{1-e}),\qquad
 p_i^-=q^{-2e-1}(1-q^{-2i}),\qquad
 p_i^0=1-p_i^+-p_i^- .
\]
The boundary factors make $p_h^+=p_0^-=0$. After $k=a/2$ steps this is the exact half-rank law of $S(Z)$, since its independent row pairs add uniform wedges. Thus \eqref{eq:positive} needs only the final mass at $h$, and \eqref{eq:negative} needs the final mass above $c$. The recursion uses $O(kh)$ rational arithmetic operations and $O(h)$ storage. The accompanying \texttt{distance.py} evaluates it for arbitrary admissible parameters.

\section{Which side of one is the likelihood on?}

\begin{proof}[Proof of Theorem~\ref{th:finite}]
Suppose first $c=h+s$ with $s\ge0$. At a nonsingular residual, Lemma~\ref{lem:like} gives the ratio $\alpha_s$ in \eqref{eq:positive}. This is strictly below one, since each factor of $\eta_{h+s}/\eta_s$ has smaller index than its counterpart in $\eta_{h+k+s}/\eta_{k+s}$.

It remains to show that already at nullity two $F_c(B)>1$; division by $\Pi\le1$ can only increase it. After fixing $h-1$ symplectic pairs, a remaining pair is chosen in dimension $2s+2$. In dimension $d=2s+2$ there are $(q^d-1)q^{d-1}$ pairs with a prescribed nonzero pairing, and $q^d+(q^d-1)q^{d-1}$ with pairing zero. Hence the density at nullity two is
\[
 \frac{\eta_{h+s}}{\eta_s}\left(1+\frac{q}{q^{2s+2}-1}\right).
\]
Writing $u=q^{-2s}\le1$, the product in front is at least $1-u/(q^2-1)$. The latter bound multiplied by $1+qu/(q^2-u)$ exceeds one, since the difference has the sign of
$q^3-q^2-q-(q-1)u\ge q^3-q^2-2q+1>0$. Lemma~\ref{lem:mono} covers every larger nullity. The mass deficit on the nonsingular event is therefore the whole total variation distance.

If $c=h-r<h$, the largest possible residual rank is $2c$. For a form of this rank the map $T$ must have kernel exactly the given radical. It induces a symplectic isomorphism on the quotient, giving $|\Sp_{2c}(q)|=q^{c(2c+1)}\eta_c$ choices by the ordered-pair count. At this rank
\[
 F_c(B)=q^{\binom b2-2cb+c(2c+1)}\eta_c
       =q^{r(2r-1)}\eta_c>1.
\]
Indeed $q\eta_c\ge q(1-1/(q^2-1))>1$ for all $q\ge2$. Monotonicity gives the same strict inequality on the entire support; the remaining points have likelihood zero. Summing the uniform mass of these points gives \eqref{eq:negative}.
\end{proof}

\section{Fourier averaging of the residual}

The unshifted Gram matrix can have a very nonuniform law when $a$ is close to $b$. Its translate by a nonsingular form has a corank law which is much closer to uniform. Let $U$ be uniform on $\Alt_b$ and $\lambda_b(j)=\Prob(\nil U=j)$.

\begin{lemma}[Corank mixing, with leading term]\label{lem:mix}
For every positive even $a,b$,
\[
 \TV(\nil S(Z),\nil U)\le \tfrac12q^{-a}\sqrt{\theta_h^{-1}-1}.
\]
More precisely, for even $j$ put
\[
 H_b(j)=\frac{(q+1)q^b-q^{b+j}-q}{q(q^2-1)(q^{b-1}-1)}.
\]
The signed error after first order satisfies
\[
 \sum_j\left|\Prob(\nil S(Z)=j)-\lambda_b(j)(1+q^{-a}H_b(j))\right|
 \le q^{-2a}\sqrt{\theta_h^{-1}-1}.
\]
\end{lemma}
\begin{proof}
Write $\nu$ for the uniform probability measure on the nonsingular orbit. A uniform alternating matrix is nonsingular with probability $\theta_h$. Indeed the orbit has size $|\mathrm{GL}_b(q)|/|\Sp_b(q)|$, giving this fraction after division by $q^m$. Replacing $H_0$ by an independent $B_0\sim\nu$ preserves the corank distribution of the translated Gram matrix, since $G_k$ is invariant under congruences. The law of $B_0+G_k$ is itself invariant. Its orbit probabilities are therefore determined by its corank.

Use the additive characters $\chi_W(B)=\psi(\sum_{i<j}W_{ij}B_{ij})$ for a fixed nontrivial additive character $\psi$ of $\F$. Then $\E\chi_W(G_k)=q^{-k\rk W}$. For one independent pair the expectation is $\E\psi(u^\T Ww)=q^{-\rk W}$, by conditioning on $u$. Multiplication over $k$ pairs gives the formula. The same holds for a Gram form with $H_1^{-1}$ in place of $J_k$ by change of basis. Put $f=\mathbf1_{\mathrm{nonsing}}/\theta_h$, the density of $\nu$ relative to uniform. Its nonconstant Fourier coefficients have squared sum $\theta_h^{-1}-1$ by Parseval. Convolution multiplies every nonconstant coefficient by at most $q^{-a}$ and every coefficient of rank at least four by at most $q^{-2a}$. Cauchy--Schwarz bounds the mean absolute error by its root mean square. This proves the first bound and the second up to the evaluation of the rank-two contribution.

Every rank-two character is indexed by a bivector $u\wedge w$ with $u,w$ independent, in exactly $|\mathrm{SL}_2(q)|=q(q^2-1)$ ways. At a matrix $B$ of nullity $j$, summing $\psi(u^\T Bw)$ over all ordered pairs gives $q^{b+j}$. The dependent pairs all give one and number $(q+1)q^b-q$. Thus the sum of the rank-two characters there is
\[
 K(B)=\frac{q^{b+j}-(q+1)q^b+q}{q(q^2-1)}.
\]
All their coefficients in $f$ are equal by transitivity. Their common value is $K(B_0)$ divided by their number, computed as $K(0)$, hence is $-1/(q^{b-1}-1)$. The inversion sign leaves $K$ unchanged. Multiplication by $q^{-a}$ gives precisely $q^{-a}H_b(j)$. Summing the absolute error over the corank classes proves the bound.
\end{proof}

\section{The profile and its first correction}

\begin{proof}[Proof of Theorem~\ref{th:limit}]
For $s\ge0$, Lemma~\ref{lem:mix} gives $\Prob(S(Z)\text{ nonsingular})=\theta_h+O_q(q^{-a})$. The product tails give $\theta_h=\theta_\infty+O_q(q^{-b})$ and, uniformly for $s\ge0$,
\[
 \alpha_s=\frac{\eta_\infty}{\eta_s}+O_q(q^{-a}+q^{-b}).
\]
For the latter bound replace each of $\eta_{h+s},\eta_{k+s},\eta_{h+k+s}$ by $\eta_\infty$, using $|\eta_j-\eta_\infty|\le q^{-2j}/(q^2-1)$ and the positive lower bound on all denominators. Insert these estimates in \eqref{eq:positive}.

For $s<0$, the uniform alternating corank has the finite law
\begin{equation}
 \lambda_b(2j)=q^{-j(2j-1)}\frac{\gamma_b}{\gamma_{2j}\eta_{h-j}}
 \quad(0\le j\le h).
 \label{eq:lambda}
\end{equation}
Choose the $2j$-dimensional radical (using the Gaussian binomial coefficient) and a nonsingular form on the quotient, and divide by $q^m$. Simplifying gives \eqref{eq:lambda}. Let $\lambda(2j)=\theta_\infty q^{-j(2j-1)}/\gamma_{2j}$. The finite probabilities are bounded by $C_q q^{-j(2j-1)}$, so dominated convergence shows $\sum\lambda=1$. The same product-tail bound and $|\gamma_b-\gamma_\infty|\le q^{-b}/(q-1)$ give
\[
 |\lambda_b(2j)-\lambda(2j)|\le C_q q^{-b}(1+q^{2j})\lambda(2j).
\]
This also holds for $j>h$, where the finite probability is zero. The right side is summable with total $O_q(q^{-b})$. Apply Lemma~\ref{lem:mix} and \eqref{eq:negative} to the event $j<-s$. The estimate holds for all thresholds simultaneously. The two tails of \eqref{eq:profile} now follow from the product and probability formulas.
\end{proof}

\begin{corollary}[First finite-size correction above the boundary]\label{cor:first}
Let $s\ge0$, $A_s=\eta_\infty/\eta_s$, $u=q^{-2s}$, and $w=q^{-a}+q^{-b}$. Uniformly over all such dimensions,
\[
 D=\mathcal D_q(s)+\frac{\theta_\infty}{q^2-1}\bigl[q(1-A_s)-uA_s\bigr]w+O_q(uw^2).
\]
If $s,a,b\to\infty$ at any relative rates, then $D\sim \theta_\infty q^{-2s}/(q^2-1)$. At $s=0$ this extends the article's square-block coefficient to unequal complementary frames.
\end{corollary}
\begin{proof}
At $j=0$, $H_b(0)=(q^b-1)/((q^2-1)(q^{b-1}-1))=q/(q^2-1)+O_q(q^{-b})$. Lemma~\ref{lem:mix} and first-order expansion of the product for $\theta_h$ give
$\Prob(S(Z)\text{ nonsingular})=\theta_\infty[1+qw/(q^2-1)]+O_q(w^2)$. Likewise $\alpha_s=A_s[1+uw/(q^2-1)]+O_q(uw^2)$. The three relevant product tails have first-order sizes $O_q(uq^{-a})$, $O_q(uq^{-b})$ and $O_q(uq^{-a-b})$, respectively. These expansions use $\prod(1-t_i)=1-\sum t_i+O((\sum t_i)^2)$ on the tails, with uniformly positive denominators. Since $1-A_s=O_q(u)$, substitution in \eqref{eq:positive} proves the claimed error. The last asymptotic uses $1-A_s=u/(q^2-1)+O_q(u^2)$.
\end{proof}

\section{Scope and further questions}

The result gives the optimal test for the full observation and a fixed-field asymptotic for every pair of nondegenerate frame types. The group experiment is a finite Fourier model isolated by the source project. Passing to it removes all hypotheses about zeros of $L$-functions. For frames with arbitrary radicals the source gives a finite likelihood. The corresponding general optimal-test thresholds and uniform limiting profile remain open. A specific obstacle is that a prescribed degenerate Gram form can be realized by a noninjective map, so the conditioning in Lemma~\ref{lem:like} must retain an extra injectivity constraint. It would also be interesting to determine the optimal test for several symplectic pairs over a finite local ring. There the reduction of a residual form to its rank over the residue field loses its higher elementary divisors.

The table gives limiting distances at adjacent ambient dimensions, rounded to six decimal places using \eqref{eq:profile}.
\begin{center}\small
\begin{tabular}{lrrrrrr}
\toprule & $s=-2$ & $s=-1$ & $s=0$ & $s=1$ & $s=2$ & $s=3$\\
\midrule
$q=2$ & 0.978652 & 0.419422 & 0.130634 & 0.034372 & 0.008702 & 0.002182\\
$q=3$ & 0.998445 & 0.639005 & 0.078878 & 0.008863 & 0.000986 & 0.000110\\
$q=5$ & 0.999933 & 0.793335 & 0.033003 & 0.001322 & 0.000053 & 0.000002\\
\bottomrule
\end{tabular}
\end{center}

A companion program checks 27 parameter choices over $\mathbb F_2$ and $\mathbb F_3$ independently by counting row pairs with a prescribed alternating Gram matrix. The finite calculations test signs, distances and Fourier estimates; the proofs above cover all prime powers. Earlier work includes the subspace counts of De Boeck and Van de Voorde \cite{anzahl}, the rank-intersection counts for alternating forms of Pott, Schmidt and Zhou \cite{pott}, and sharp convergence estimates for classical finite-field rank laws by Fulman and Goldstein \cite{stein}. The contributions claimed here are the optimal residual-rank tests away from complementary dimension, the uniform profile for the compression experiment, and the uniform first-order description of the translated Gram corank. Priority beyond the sources compared remains a bibliographical question.

\begin{thebibliography}{9}
\bibitem{proveit} ProveIt Contributors, \emph{Beyond the Riemann hypothesis for Dirichlet $L$-functions: consequences, transfers and limits}. ProveIt, snapshot \texttt{35e2a5a63492} (2026). \href{\src}{Section on spanning, frames and matrix statistics}. See \texttt{spn:prop:residual}, \texttt{spn:thm:balanced}, \texttt{spn:thm:square-block} and \texttt{spn:q:symplectic}.
\bibitem{anzahl} M. De Boeck and G. Van de Voorde, \emph{Anzahl theorems for trivially intersecting subspaces generating a non-singular subspace I: symplectic and hermitian forms}, Linear Algebra Appl. \textbf{699} (2024). \href{https://doi.org/10.1016/j.laa.2024.07.004}{doi:10.1016/j.laa.2024.07.004}. See Theorem 4.9, \href{https://arxiv.org/abs/2407.07486}{arXiv:2407.07486}.
\bibitem{pott} A. Pott, K.-U. Schmidt and Y. Zhou, \emph{Pairs of quadratic forms over finite fields}, Electron. J. Combin. \textbf{23} (2016), P2.8. \href{https://doi.org/10.37236/4072}{doi:10.37236/4072}.
\bibitem{stein} J. Fulman and L. Goldstein, \emph{Stein's method and the rank distribution of random matrices over finite fields}, Ann. Probab. \textbf{43} (2015), 1274--1314. \href{https://doi.org/10.1214/13-AOP889}{doi:10.1214/13-AOP889}.
\end{thebibliography}
\end{document}
